\documentclass[margin,line]{res}


\oddsidemargin -.5in
\evensidemargin -.5in
\textwidth=6.0in
\itemsep=0in
\parsep=0in
% if using pdflatex:
%\setlength{\pdfpagewidth}{\paperwidth}
%\setlength{\pdfpageheight}{\paperheight} 

\newenvironment{list1}{
  \begin{list}{\ding{113}}{%
      \setlength{\itemsep}{0in}
      \setlength{\parsep}{0in} \setlength{\parskip}{0in}
      \setlength{\topsep}{0in} \setlength{\partopsep}{0in} 
      \setlength{\leftmargin}{0.17in}}}{\end{list}}
\newenvironment{list2}{
  \begin{list}{$\bullet$}{%
      \setlength{\itemsep}{-0in}
      \setlength{\parsep}{0in} \setlength{\parskip}{0in}
      \setlength{\topsep}{0in} \setlength{\partopsep}{0in} 
      \setlength{\leftmargin}{0.2in}}}{\end{list}}
%\setlength{\textheight}{9.5in} % increase text height to fit on 1-page 
\usepackage{color,hyperref}
\usepackage{eurosym}
\usepackage{multicol}

\usepackage{enumitem}

\begin{document}

\name{ILARIA COLAZZO \vspace*{.1in}}

\begin{resume}
\small
\section{\sc Contact Information}
\vspace{.05in}
\begin{tabular}{@{}p{4in}p{8in}}
University of Exeter  &  \href{mailto:I.Colazzo@exeter.ac.uk}{I.Colazzo@exeter.ac.uk} \\       
Department of Mathematics& \href{mailto:ilariacolazzo@gmail.com}{ilariacolazzo@gmail.com}\\
Exeter, UK &  \href{https://www.ilariacolazzo.info}{www.ilariacolazzo.info} \\
\end{tabular}

%\vspace{-0.1in}
\section{\sc Current \\Position}
{\bf  University of Exeter}, Exeter, UK\\ %{\bf October 1, 2019}\\
{\em Postdoctoral Research fellow} 
{}\\
\vspace*{-.1in}

%\vspace{-0.2in}
\section{\sc Research\\ Interests} 

Algebra,
Set-theoretical solutions of the Yang-Baxter equation,
Skew braces, trusses and generalisations,
Regular subgroups of the holomorph,
Hopf-Galois extensions,
Set-theoretical solutions of the pentagon equation

%\vspace{-0.1in}
\section{\sc Research Experience}

{\bf University of Exeter}, Exeter, UK, \hfill {\bf 01/06/2021 -- Present}

\vspace{-.3cm}
{\em Postdoctoral Researcher in Mathematics} \\
{EPSRC project: Hopf-Galois Theory and Skew Braces.  }\\
{PI: Prof. Nigel Byott}
\vspace{3pt}
%\begin{list2}
%	\item  Studied the algebraic structure of a semi-truss. Deepen the connection between set-theoretic solutions associated to a semi-truss and set-theoretic solutions associated to semi-braces. Classified the involutive solutions of the Pentagon equation.
%\end{list2}


\vspace{-0.1in}

{\bf Vrije Universiteit Brussel}, Brussels, Belgium, \hfill {\bf 01/10/2019 -- 28/02/2021}

\vspace{-.3cm}
{\em Postdoctoral Researcher in Mathematics} \\
%{Research project: ”Yang-Baxter equation and related algebric structures”.  }\\
{Scientific advisor: Prof. Eric Jespers}
\vspace{3pt}
%\begin{list2}
%	\item  Studied the algebraic structure of a semi-truss. Deepen the connection between set-theoretic solutions associated to a semi-truss and set-theoretic solutions associated to semi-braces. Classified the involutive solutions of the Pentagon equation.
%\end{list2}

\vspace{-0.1in}
{\bf University of Salento}, Lecce, Italy\hfill {\bf 01/12/2017 -- 30/11/2018}

\vspace{-.3cm}
{\em Postdoctoral Researcher in Mathematics} \\
{Research project: ”Yang-Baxter equation and related algebraic structures”.  }\\
{Scientific advisor: Prof. Francesco Catino}
\vspace{3pt}
%\begin{list2}
%	\item Introduced a novel construction technique for set-theoretical solutions of the Yang-Baxter equation that to obtain new classes of involutive and idempotent solutions. This technique led to characterizing completely the solutions associated with a finite semi-brace.
%\end{list2}

\vspace{-0.1in}
{\bf Vrije Universiteit Brussel}, Brussels, Belgium \hfill{\bf 01/09/2018 -- 30/11/2018}

\vspace{-.3cm} 
{Postdoctoral Researcher Visiting (host Prof. E. Jespers)}
% \begin{list2}
	% 	\item Established collaborations with E. Jespers, A. Van Antwerpen to obtain a description of  the algebraic structure of semi-trusses.
	% \end{list2}

\vspace{-0.1in}
{\bf University of Salento and University of Basilicata}, Lecce, Italy \hfill{\bf 03/03/2014 -- 20/07/2017}


\vspace{-.3cm}
{\em PhD Student}
%\begin{list2}
%	\item Presented new construction techniques (i.e., the  Hochschild product and the asymmetric product) for regular subgroups of an affine group using their relationship with braces over a field.
%	Advanced the state-of-the-art of the algebraic structures related to set-theoretical solutions of the Yang-Baxter equation by introducing semi-braces.
%\end{list2}

\vspace{-0.1in}
{\bf University of Warsaw}, Warsaw, Poland \hfill{\bf 15/11/2016 -- 20/12/2016}

\vspace{-.3cm}
{\em Visiting PhD Student} (host Prof. J. Okniński)
% \begin{list2}
	%\item Collaborated with J. Okniński to study solutions of the Yang-Baxter equation which are not necessarily involutive.
	%\end{list2}
	
\vspace{8pt}
\section{\sc Education}
{\bf University of Salento}, Lecce, Italy \hfill {\bf  20/07/2017}\\
{\em Department of Mathematics and Physics ``Ennio De Giorgi''} 
%\vspace*{-.1in}
\begin{list1}
\item[] Ph.D., Mathematics and Informatics
\begin{list2}
\vspace*{.05in}
\item[] Dissertation Title: ``Left Semi-Braces and the Yang-Baxter equation''. 
\item[] Advisor: Prof. F. Catino
\item[] Final mark: cum Laude
\end{list2}
\end{list1}

\vspace{-0.1in}
{\bf University of Salento and University of Basilicata}, Lecce, Italy \hfill {\bf 25/10/2012} \\
{\em Department of Mathematics and Physics} 
%\vspace*{-.1in}
\begin{list1}
\item[] M.S., Mathematics
\begin{list2}
%	\vspace*{.05in}
%	\item[] Dissertation Title: ``Clifford algebras''. 
%	\item[] Supervisor: Prof. F. Catino
	\item[] Final mark: 110/110 cum Laude
\end{list2}
\end{list1}

\vspace{-0.1in}
{\bf University of Salento}, Lecce, Italy\hfill {\bf  15/06/2009}\\
{\em Scuola Normale Superiore ISUFI}\\
{\em Nanoscience and Grid-Computing}
\begin{list1}
	\item[] Licence (it: Licenza Scuola Superiore ISUFI Pre-Laurea I livello)
%	\begin{list2}
%		\item[] Dissertation Title: ``Development of geometry manipulation tools for atomistic simulations''.
%		\item[] Supervisor: Prof. M. De Vittorio
%	\end{list2}
\end{list1}

\vspace{-0.1in}
{\bf University of Salento}, Lecce, Italy\hfill {\bf  18/12/2008} \\
{\em Department of Mathematics and Physics} 
\begin{list1}
\item[] B.S., Mathematics 
\begin{list2}
%	\vspace*{.05in}
%	\item[] Dissertation Title: ``On the Galois group of a polynomial of degree five".
%	\item[] Supervisor: Prof. M. M. Miccoli
	\item[] Final mark: 110/110 cum Laude
\end{list2}
\end{list1}





%\section{\sc Publications}
%\vspace{-0.1in}
%\begin{enumerate}[{label=[\arabic{*}]}, leftmargin=*]
%	\item {\bf I. Colazzo}, E. Jespers, A. Van Antwerpen, C. Verwimp, \href{https://doi.org/10.1016/j.jalgebra.2022.07.019}{\emph{Left non-degenerate set-theoretic solutions of the Yang-Baxter equation and semitrusses}}, J. Algebra (2022).
%	\item {\bf I. Colazzo}, E. Jespers,  Ł. Kubat. \href{https://doi.org/10.1007/s00220-020-03862-6}{\emph{Set-Theoretic Solutions of the Pentagon Equation}}, Commun. Math. Phys. (2020).
%	\item {\bf I. Colazzo}, A. Van Antwerpen, \href{https://doi.org/10.1016/j.jpaa.2020.106467}{\emph{The algebraic structure of left semi-trusses}}, J. Pure Appl. Algebra, 225 (2021) n.2, 106467.
%	\item F. Catino, {\bf I. Colazzo}, P. Stefanelli, \href{https://doi.org/10.1007/s00009-020-1483-y}{\emph{The Matched Product of the Solutions to the Yang–Baxter Equation of Finite Order}}, Mediterr. J. Math., 17 (2020), no. 2, Art. 58.
%	\item F. Catino, {\bf I. Colazzo}, P. Stefanelli, \href{https://doi.org/10.1007/s00009-020-1483-y}{\emph{The Matched Product of the Solutions to the Yang–Baxter Equation of Finite Order}}, Mediterr. J. Math., 17 (2020), no. 2, Art. 58, .
%	\item F. Catino, {\bf I. Colazzo}, P. Stefanelli, \href{https://doi.org/10.1016/j.jpaa.2019.07.012}{\emph{The matched product of solutions to the Yang-Baxter equation}}, J. Pure Appl. Algebra 224 (2020) n. 3, 1173--1194.
%	\item F. Catino, {\bf I. Colazzo}, P. Stefanelli, \href{https://doi.org/10.1142/S0219498819500336}{\emph{Skew left braces with non-trivial annihilator}}, J. Algebra Appl. 18 (2019) n.2, 1950033, 23 pp.
%	\item F. Catino, {\bf I. Colazzo}, P. Stefanelli, \href{https://doi.org/10.1016/j.jalgebra.2017.03.035}{\emph{Semi-braces and the Yang–Baxter equation}},	J. Algebra, \textbf{483}, (2017), 163--187.
%	\item F. Catino, {\bf I. Colazzo}, P. Stefanelli, \href{http://dx.doi.org/10.1016/j.jalgebra.2016.01.038}{\emph{Regular subgroups of the affine group and asymmetric product of radical braces}}, J. Algebra, \textbf{455}, (2016), 164--182.
%	\item F. Catino, {\bf I. Colazzo}, P. Stefanelli, \href{http://dx.doi.org/10.1017/S000497271400077X}{\emph{On regular subgroups of the affine group}}, Bull. Aust. Math. Soc., \textbf{91} 		
%	(2015), 76--85.
%\end{enumerate} 
%
%\section{\sc Preprints and others}
%\vspace{-0.1in}
%\begin{enumerate}[{label=[\arabic{*}]}, leftmargin=*]
%	\item {\bf I. Colazzo}, M. Ferrara, M. Trombetti, \href{https://arxiv.org/abs/2210.08598}{\em On derived-indecomposable of the Yang-Baxter equation}, \\arXiv:2210.08598.
%	\item {\bf I. Colazzo}, \href{https://www.aracneeditrice.eu/it/pubblicazioni/algebra-for-cryptography-riccardo-aragona-norberto-gavioli-filippo-mignosi-9791259943286.html}{\em Braces: between regular subgroups and solutions of the Yang-Baxter equation}, Algebra for Cryptography, Collectio Ciphrarum, (Extended abstract), ISBN 9791259943286 (2021).
%\end{enumerate}

ations}

\nocite{*}

\newrefcontext[labelprefix=P]
\printbibliography[type=unpublished,title={Preprints},resetnumbers=true,heading=subbibnumbered]

\newrefcontext[labelprefix=J]
\printbibliography[type=article,title={Journal papers},resetnumbers=true,heading=subbibnumbered]

\newrefcontext[labelprefix=C]
\printbibliography[type=inproceedings,title={Refereed proceedings},resetnumbers=true,heading=subbibnumbered]

\newrefcontext[labelprefix=O]
\printbibliography[title={Other publications},resetnumbers=true,filter=other,heading=subbibnumbered]

\section{Talks}

\newrefcontext[labelprefix=I]
\printbibliography[title={Invited talks},resetnumbers=true,filter=invited,heading=subbibnumbered]

\newrefcontext[labelprefix=T]
\printbibliography[title={Confe
 
\vspace{-0.1in}
\section{\sc  Invited\\ Talks}

{\sc Invited talks at conferences and workshops}

\vspace{-0.1in}
\emph{YB-semitrusses and left non-degenerate solutions to the Yang-Baxter equation}:
Conference Hopf algebras and Galois module theory (Omaha/Online) 30 May -- 03 June 2022.

\vspace{-0.1in}
\emph{YB-semitrusses with associated bijective solutions}:
Braces in Bracelet Bay, Swansea University (Online), 04--06 June 2022

\vspace{-0.1in}
\emph{Skew braces and solutions of the Yang-Baxter equation}:
Conference  Hopf algebras and Galois module theory (Online), 24--27 May 2021

\vspace{-0.1in}
\emph{Braces: between regular subgroups and solutions of the Yang-Baxter equation}: Workshop \emph{Algebra for Cryptography} (L'Aquila), 10-11 October 2019.


\vspace{-0.1in}
\emph{The matched product of the solutions of the Yang-Baxter equation}: Meeting \emph{Noncommutative and non-associative structures, braces and applications} (Malta), 11-15 March 2018.


\vspace{-0.1in}
{\sc Seminar Talks}

\vspace{-0.1in}
\emph{Retractability problem for quasi-linear cycle sets}, Seminar Algebra, University of Warsaw (Poland), 20 October 2022.

\vspace{-0.1in}
\emph{Set-theoretic solution of the Pentagon Equation: the involutive case}, QA Seminar, Vrije Universiteit Brussel (Belgium), 18 May 2022.

\vspace{-0.1in}
\emph{Set-theoretic solutions of the Yang-Baxter equation and related associative structure}, University of Padua (Italy), 20 April, 2022.

\vspace{-0.1in}
\emph{Bijective set-theoretic solutions of the Pentagon Equation},  Al@Bicocca take-away, University of Milan – Bicocca, Milan (Italy), 12 November 2021.

\vspace{-0.1in}
\emph{Set-theoretic solutions of the Yang-Baxter equation, skew braces and regular subgroups}, University of Exeter (UK), 11 November 2021.

\vspace{-0.1in}
\emph{Regular subgroups and left semi-braces}, ALGB Seminar, Vrije Universiteit Brussel (Belgium),  31 October 2018.

\vspace{-0.1in}
\emph{Regular subgroups of the affine group, Seminar Algebra},  Seminar Algebra, University of Warsaw (Poland) 24 November 2016.


\vspace{-0.1in}


%\emph{Regular subgroups of an affine group}, Meeting \emph{Group Theory in Florence - A meeting in honour of Guido Zappa}: Florence (Italy), June 16-17, 2016.
%
%\emph{The Asymmetric Product of radical braces}, International Conference \emph{Advances in Group Theory and Applications 2015}, Porto Cesareo (Lecce, Italy), June 16-19, 2015.

%\section{\sc Posters}
%\emph{The Asymmetric Product: a new construction of radical $F$-braces}, Summer school \emph{Advances in Group Theory and Applications 2016 - The School}: Vietri sul Mare (Salerno, Italy), June 6-10, 2016.
%
%\emph{New construction of radical $F$-braces: The Hochschild product}, International Conference \emph{Naples 2015 -- Conference in 
%		Group Theory and its applications}: Naples (Italy), October 7-8, 2015. 2015.
%
%\section{\sc Seminar \\Talks}

%\emph{Sui brace di ordine $p^2q$}, Algebra Reading Seminar, University of Salento (31/01/2018)\\
%\emph{Hopf braces}, Algebra Reading Seminar, University of Salento(31/10/2017)\\
%\emph{Regular subgroups of the affine group}, Seminar Algebra, University of Warsaw (24/11/2016)\\
%\emph{Il prodotto asimmetrico di brace}, Algebra Reading Seminar, University of Salento (29/01/2016)\\
%\emph{Algebre di Hopf quasi-cocommutative}, Algebra Reading Seminar, University of Salento (29/10/2015)

%\vspace{0.3cm}
%\section{\sc Summer\\Schools}
%\emph{Advances in Group Theory and Applications 2016 - The School}, Vietri sul Mare (Salerno, Italy), June 6-10, 2016, organized by F. Catino, M. De Falco, F. de Giovanni, C. Musella, courses held by C. Casolo, \'{A}. del R\'{i}o, A. Facchini, A. J. Macintyre.
%
%\emph{Bicocca Ph.D. School on Representation Theory 2014}, Università di Milano- Bicocca (Italy), June 16-18, 2014, organized by M. Avitabile, F. Dalla Volta, L. Di Martino, A. Previtali, P. Spiga and Th. Weigel, courses held by M. Isaacs and G. Navarro.
%
%\emph{Advances in Group Theory and Applications 2014 - The School}, Porto Cesareo (Lecce, Italy), June 3-6, 2014, organized by F. Catino and F. de Giovanni, courses held by E. Aljadeff, A. Ballester-Bolinches, F. de Giovanni and H. Laue.
%
%
%\vspace{0.4cm}
\vspace{-0.1in}
\section{\sc Conferences \\ Organised}


%\vspace{3pt}
\begin{list2}
	
	\item \emph{Oberwolfach mini-workshop (2309a)}  \hfill{\bf	26/02--03/03/2023}\\
	\emph{Skew Braces and the Yang–Baxter Equation}\\
	Oberwolfach (Germany), member of the scientific committee\\ 
	(with T. Brzezinski, A. Doikou, and L. Vendramin)
	
	\item \emph{The algebra of the Yang–Baxter equation} \hfill{\bf 10--16/07/2019} \\
	Bedlewo (Poland), member of the organising \\
	(with Ł. Kubat, V. Lebed, J. Okniński, A. Van Antwerpen, L. Vendramin, and M. Wiertel)
	
	\item  \emph{Conference Young Researchers Algebra Conference 2019}\hfill{\bf16--19/09/2019} \\
	Naples (Italy), member of the scientific committee \\
	(with M. Brescia, M. Ferrara, P. Stefanelli and M. Trombetti)
	
	\item \emph{Advances in Group Theory and Applications 2019}\hfill{\bf 25--28/06/2019}\\
	  Lecce (Italy), member of the organising committee \\
	 (with M. Brescia, M. Ferrara, M. M. Miccoli, P. Stefanelli and M. Trombetti)
	
	\item \emph{Advances in Group Theory and Applications 2017} \hfill{\bf 05--08/09/2017} \\
	Lecce (Italy), member of the organising committee\\
	(with M. M. Miccoli, P. Stefanelli and M. Trombetti)
	
	\item \emph{Advances in Group Theory and Applications 2015}\hfill{\bf 16–19/06/2015}\\
	Porto Cesareo, Lecce (Italy), member of the organising committee\\ 
	(with M. M. Miccoli, P. Stefanelli)
	
	\item \emph{Algebra Reading Seminar} \hfill {\bf 2015–2019}\\
	Co-organiser of the Department seminar (with F. Catino)\\
	University of Salento, Department of Mathematics and Physics “Ennio De Giorgi”
 	
% 		\vspace{.3cm} 
% 	\item {\em Representative Ph.D. Student} (University of Salento)\hfill {\bf 2014 -- 2016}
% 	 	\item {\em Representative Student} (University of Salento)\hfill{\bf 2010 -- 2012}
\end{list2}



\section{\sc Teaching experience}

{\bf MAGIC} (UK)\\
{\em a collaboration of 22 universities, coordinated by the University of Exeter.}

%\vspace{-0.04in}
%\vspace{.3cm}
\begin{list2}
\item {\em Lecturer}
\item[]	Duties included preparing and delivering lectures, preparing study material and exercises, marking homework and the final assignment.
\begin{list2}
	\item[$\circ$] MAGIC111 \hfill{A.Y. 2022–23}
	\item[] Introduction to set-theoretic solutions to the Yang-Baxter equation and skew braces
	\item[] Course for PhD students
%	\item[] (for a total of 10 hours)
\end{list2}
\end{list2}

%\vspace{-0.1in}
{\bf University of Exeter}, Exeter, UK
%{\em Department of Mathematics}
%\vspace{-0.04in}
\begin{list2}
	\item {\em Tutor} 
	\item[] Duties included preparing lectures for tutorial classes, fielding of student inquiries, and marking homework.
	\begin{list2}
		\item[$\circ$] \emph{MTH2010 – Groups, Rings and Fields} (Course Prof M. Saidi)  \hfill{ A.Y. 2022–23}
%		\item[] (for a total of 10 hours)
		\item[$\circ$] \emph{MTH1001 – Mathematical structures} (Course Dr G. Marasingha) \hfill{ A.Y. 2021–22}
%		\item[](for a total of 10 hours)
\end{list2}
\end{list2}

%\vspace{-0.1in}
{\bf University of Salento}, Lecce, Italy
%{\em Department of Mathematics and Physics ``Ennio De Giorgi''}

%\vspace{-0.04in}
%\vspace{.3cm}
\begin{list2}
	\item {\em Instructor}
	\item[] Duties included preparing lectures, fielding of student inquiries, and office hours.
%	\vspace{.1cm}
	\begin{list2}
		\item[$\circ$]  Mathematics preliminary  course \hfill {A.Y. 2017--18}\\for bachelor's students in Mathematics and in Physics 
	\end{list2}
%	\vspace{.3cm}
	\item \emph{Teaching assistant} 
	\item[] Duties included preparing lectures for tutorial classes for recitation classes and office hours.
	\vspace{.1cm}
	\begin{list2}
		\item[$\circ$] Group Theory (Course Professor F. Catino) \hfill {A.Y. 2017--18}
		\item[$\circ$] Group Theory (Course Professor F. Catino)\hfill {A.Y. 2016--17}
		\item[$\circ$] Group Theory (Course Professor F. Catino)\hfill {A.Y. 2015--16}
		\item[$\circ$] Algebra II (Course Professor M. M. Miccoli)\hfill {A.Y. 2015--16}
	\end{list2} 
%	\vspace{.3cm}
	\item {\em Tutor} 
	\item[] Duties included preparing lectures for tutorial classes, fielding of student inquiries, and office hours.
	\vspace{.1cm}
	\begin{list2}
		\item[$\circ$] Algebra I and II (Course Prof F. Catino and Prof M. M. Miccoli) \hfill {A.Y. 2018--19}
		\item[$\circ$] Algebra I (Course Prof F. Catino)\hfill {A.Y. 2016--17}
	\end{list2}
%	\vspace{.2cm}
	\item {\em Algebra, Exam committee member} (it Cultore della materia) \hfill {July 2016 -- July 2019}
	\item[] Duties included grading final exams.
\end{list2} 




\section{\sc Refereeing \\activity}
Annales Henri Poincare, Journal Publicacions Matemàtiques, Journal of the London Mathematical Society (JLMS), International Mathematics Research Notices, Journal of Algebra and its Applications, Communications in Algebra, Rendiconti del Circolo Matematico di Palermo Series 2.

%
\section{\sc Skills} 
{\bf Programming Language:} GAP, Python.

\vspace{-.3cm}
{\bf Language:} Italian (native), English (fluent), French (waystage).


%\section{\sc References} 
%\begin{tabular}{ll}
%	Prof. Eric Jespers& Prof. Francesco Catino\\
%	Department of Mathematics&Department of Mathematics and Physics\\
%	Vrije Universiteit Brussel& University of Salento\\
%	Pleinlaan 2, B-1050 Brussels, Belgium & Monteroni, Lecce, 73100, Italy\\
%	\href{mailto:eric.jespers@vub.be}{eric.jespers@vub.be}&\href{mailto:francesco.catino@unisalento.it}{francesco.catino@unisalento.it}\\
%	\mbox{}&\\
%	Prof. Francesco de Giovanni &\\
%	Dipartimento di Matematica e Applicazioni&\\ "Renato Caccioppoli"&\\
%	Università degli studi di Napoli
%	Federico II&\\
%	Via Cinta - Complesso Monte S. Angelo, 29&\\
%	\href{mailto:degiovan@unina.it}{degiovan@unina.it}
%\end{tabular}













\end{resume}
\end{document}




